\section{Sorts and Notation}

\subsection{Sorts}

\begin{description}
  \item[$\Name$] A path component: a finite sequence of bytes containing neither
        \code{/} nor NUL.  On macOS it is always valid UTF-8, and is read as the
        sequence of Unicode characters it encodes; on Linux it may be any such bytes
        (Law~\ref{law:wire} says how the client is given one that is not UTF-8).
  \item[$\Path$] An absolute path, written as the sequence of its components
        $p=\Seq{n_1,\dots,n_k}$; the root is $\Seq{}$.  $\Comps(p)$ is that
        sequence, $\Parent(p)=\Seq{n_1,\dots,n_{k-1}}$, and $p\cdot n$ appends a
        component.
  \item[$\Type$] $\{\Dir,\File\}$, the type of the final component.
  \item[$\Rule$] A parsed pattern in the gitignore dialect, possibly negated and
        possibly restricted to directories.
  \item[$\RuleSet$] A finite sequence of rules.  Two sets are used: $\Deny$
        (hard block, no negation) and $\Hide$ (cosmetic, negation allowed).
  \item[$\Entry$] A directory entry as shown in a listing.
  \item[$\Row$] An entry as the frontend sorts it: name, type, size, and
        modification time.
  \item[$\Conn$] A TCP connection to \texttt{fsb}, given by the server's and the
        client's address and port.
  \item[$\Uid$] A user, as the system's numeric user ID.
\end{description}

\subsection{Environment}

Let $\Roots$ be the finite set of served roots, each already resolved to its
real location.  The filesystem is treated as a partial function from paths to
files and directories with symbolic links; $\Real(p)$ is the canonical real
location of $p$ after all symbolic links and aliases (such as the macOS
firmlink prefix \code{/System/Volumes/Data}) are resolved, defined when $p$
exists.  $\Canon$ reduces aliases of the same location to one spelling.

For the session: $\tau$ is the launch token of this start, $u\in\Uid$ the user
\texttt{fsb} runs as, and $\Used\in\Bool$ whether the token has been exchanged
(initially false).  $\Owner:\Conn\rightharpoonup\Uid$ is the user owning the
client end of a connection.  On Linux it is defined when the connection table
lists that end with a user \texttt{fsb}'s user namespace can name; on macOS it is
never consulted.

\subsection{Operations}

The operations analyzed are, on the server:
\[
  \Open,\ \List,\ \Search,\ \text{and the endpoints built on them,}
\]
namely file download, head, metadata, image and PDF preview, and archive
listing; the launch exchange $\Exchange$; and on the client the pure functions $\Cmp$ and $\Sort$ (row order),
$\Filter$, $\ToHash$ and $\FromHash$, $\Goto$, and $\Link$.

\subsection{Conventions}

A relation $x\approx y$ between names means \emph{the filesystem treats $x$ and
$y$ as the same name}; it is defined in Section~\ref{sec:names}.  A statement
``$\Open(p)$ fails as $\NotFound$'' means that the operation fails with an error
that the HTTP layer answers with an identical 404.
